\documentclass[10pt,a4paper]{amsart}
\usepackage[margin=19mm]{geometry}
\usepackage{amsmath,amssymb,mathtools}
\usepackage[scr=rsfs,cal=euler]{mathalfa}
\usepackage{xfrac}
\usepackage[unicode,hidelinks,bookmarks=false]{hyperref}
\newcommand{\ie}{i.e.\ }
\newcommand{\cf}{cf.\ }
\newcommand{\resp}{respectively\ }
\newcommand{\Subsection}{\subsection}
\providecommand{\nobreakdash}{\nobreak\hbox{-}}
\setlength{\emergencystretch}{2em}
\multlinegap0pt
\usepackage{bm}
\usepackage{bbm}
\usepackage{dsfont}
\usepackage{xspace}
\usepackage{cancel}
\usepackage{mathtools}
\usepackage{amsthm}
\makeatletter
\@ifundefined{theorem}{\newtheorem{theorem}{Theorem}}{}
\@ifundefined{lemma}{\newtheorem{lemma}[theorem]{Lemma}}{}
\makeatother
\title[Detected Seifert surfaces and intervals of left-orderable su]{Detected Seifert surfaces and intervals of left-orderable surgeries}
\author{Yi Wang}
\date{Source: arXiv:2509.08127v1 (CC BY 4.0)}
\begin{document}
\maketitle

\noindent\emph{Source and licence.} Detected Seifert surfaces and intervals of left-orderable surgeries. Version arXiv:2509.08127v1 (CC BY 4.0), distributed under the Creative Commons Attribution 4.0 International licence, as recorded in the arXiv licence metadata for this exact version and shown on the abstract page https://arxiv.org/abs/2509.08127v1. Full rights notice: licenses/arxiv-2509.08127-CC-BY-4.0.txt.\par\smallskip

\noindent\textit{Editorial scope.} Counted unit: the Lemma whose source label is \texttt{lma:limiting} in the article (source0.tex lines 554--576, source label \texttt{lma:limiting}), delivered together with the article's own complete original proof of it (source0.tex lines 578--611, one \texttt{proof} environment of 34 lines). No mathematics is written, repaired or re-derived: statement and proof are the article's own text line for line. Required context retained uncounted: none. Every definition, display and label of those lines is retained unchanged. Omitted from this package and not redistributed: 1--553, 577--577, 612--789. Nothing inside a retained excerpt is omitted or altered: every retained body is delivered exactly as the article's lines are written, so no omission receipt of any kind accompanies this package. Citations: the retained excerpts carry no citation command at all, so no bibliography entry of the article is transcribed and no reference is redistributed. Reference adaptation: none was needed and none was made; every reference inside the delivered unit resolves inside it, so no unresolved reference or citation is produced. Printed slips kept literal: no printed word, symbol, hypothesis or number of the counted statement or of its proof was altered. Layout and class: the article's own class is \texttt{article} with packages including inputenc, babel, blindtext, amssymb, amsthm, amsmath, tikz-cd, setspace, thmtools, thm-restate, hyperref, cleveref, geometry, graphicx; the wrapper is the lane's pinned \texttt{amsart} editorial wrapper at 10pt on A4, which restates the theorem-like environments the retained text needs and adds the article's own macro definitions that the retained text needs (none), with extra wrapper packages (bm, bbm, dsfont, xspace, cancel, mathtools). The delivered numbering is the unit's own. Assets: no retained span contains a figure environment, a graphics-inclusion command, a PDF-page inclusion, a table or a TikZ picture; no image file is copied into this package, the package declares no asset dependency and no image file is redistributed with it. Licence: version arXiv:2509.08127v1 of this article is distributed under the Creative Commons Attribution 4.0 International licence, as recorded in the arXiv licence metadata for this exact version and shown on the abstract page https://arxiv.org/abs/2509.08127v1. A full rights notice is delivered at licenses/arxiv-2509.08127-CC-BY-4.0.txt. Characters: every retained excerpt is pure ASCII, so the delivered engine, XeLaTeX with its default Latin Modern text font, reports no missing character.
\begin{lemma}\label{lma:limiting}
Let $\rho_n: F_2 \to SL_2(\mathbb{R})$ be the representation 
\begin{equation}
	\rho_n(a) = \begin{pmatrix}-1 & 1 \\ 0 & -1\end{pmatrix} \ \ \ \ \ \rho_n(b) = \begin{pmatrix}2n+1 & n \\ 2 & 1\end{pmatrix}
\end{equation}
with $\chi_n = (-2, 2n+2, -2n)$ its character. Let 
\begin{equation}
	m_1 = a^{n+1}bab \ \ \ \ \ m_2 = a^{n+1}ba \ \ \ \ \ \ell_1 = b^{-1}ab \ \ \ \ \ \ell_2 = b^{-1}aba
\end{equation}
and define a curve
\begin{equation}
	C_n = \{\text{tr}(m_1) = \text{tr}(m_2), \text{tr}(\ell_1) = \text{tr}(\ell_2), \text{tr}(m_1\ell_1) = \text{tr}(m_2\ell_2)\} \subset \mathbb{C}^3
\end{equation}
Then the following is true.
\begin{enumerate}
	\item $\rho_n(m_1)$ and $\rho_n(m_2)$ are conjugate parabolic matrices.
	\item $\chi_n \in C_n$. 
	\item $\chi_n$ is a smooth point on $C_n$.
	\item $\text{tr}(m_1)$ is a local coordinate on $C_n$.
	\item The trace of the longitude is non-constant at $\chi_n$.
	\item $\chi_n$ corresponds to an ideal point $X(M_{-3, 3, 2n+1})$ detecting the Seifert surface, and this ideal point lies on a component of $X(M)$ on which the trace of the longitude is nonconstant.
\end{enumerate}
\end{lemma}

\begin{proof}
A computation yields
\begin{equation}
	\rho_n(m_1) = (-1)^n\begin{pmatrix}-2n-1 & -n \\ 4n & 2n-1\end{pmatrix} \ \ \ \ \ \rho_n(m_2) = (-1)^n\begin{pmatrix}-1 & 0 \\ 2 & -1\end{pmatrix}
\end{equation}
Let 
\begin{equation}
	P_1 = \begin{pmatrix}1 & 0 \\ 2 & 1 \end{pmatrix} \ \ \ \ \ P_2 = \begin{pmatrix} 1 & -1 \\ 1 & 0 \end{pmatrix}
\end{equation}
\begin{equation}
	P_1\rho_n(m_1)P_1^{-1} = (-1)^n\begin{pmatrix}-1 & -n \\ 0 & -1\end{pmatrix} \ \ \ \ \ P_2\rho_n(m_2)P_2^{-1} = (-1)^n\begin{pmatrix}-1 & -2 \\ 0 & -1\end{pmatrix} 
\end{equation}
These matrices will always have the same sign of the upper-right entry, proving \textbf{(1)}. We also have 
\begin{equation}
	\rho_n(\ell_1) = \begin{pmatrix}1 & 1 \\ -4 & -3\end{pmatrix} \ \ \ \ \ \rho_n(\ell_2) = \begin{pmatrix}-1 & 0 \\ 4 & -1\end{pmatrix}
\end{equation}
Notice that $\rho_n(\ell_1)$ and $\rho_n(\ell_2)$ aren't conjugate. Finally,
\begin{equation}
	\rho_n(m_1\ell_1) = (-1)^n\begin{pmatrix}2n-1 & n-1 \\ -4n & 3-2n\end{pmatrix} \ \ \ \ \ \rho_n(m_2\ell_2) = (-1)^n\begin{pmatrix}1 & 0 \\ -6 & 1\end{pmatrix}
\end{equation}
Then $\rho_n$ satisfies the equations $\text{tr}(m_1) = \text{tr}(m_2), \text{tr}(\ell_1) = \text{tr}(\ell_2), \text{tr}(m_1\ell_1) = \text{tr}(m_2\ell_2)$, meaning that $\chi_n = (-2, 2n+2, -2n) \in C_n$. The Jacobian at $\chi_n$ is given by 
\begin{equation}
	\begin{pmatrix}(-1)^{n+1}\frac{4n^4 + 4n^3 - 7n^2 - 4n}{3} & (-1)^{n+1}(2n^2 - n - 1) & (-1)^{n+1}(2n^2 + n - 2) \\ (2n+1)^2 & 4 & 4 \\ (-1)^{n+1}\frac{-4n^4 + 4n^3 + 43n^2 + 26n + 3}{3} & (-1)^{n+1}(-2n^2 + 5n + 9) & (-1)^{n+1}(-2n^2 + 3n + 14) \end{pmatrix}
\end{equation}
The determinant of this matrix is zero, and the top left $2 \times 2$ minor is nonzero, so the rank is 2. This confirms that $\chi_n$ is a smooth point on $C_n$, proving \textbf{(3)}. In order to show that $\text{tr}(m_2)$, we show that the gradient of $\text{tr}(m_2)$ is not in the span of the gradient. We compute that the gradient is 
\begin{equation}
	\nabla \text{tr}(m_2)|_{\chi_n} = (-1)^{n+1}\left(\frac{2n(n+1)(n+2)}{3}, n+1, n+2\right)
\end{equation}
which can be seen not to be in the span of the Jacobian of $C_n$ at $\chi_n$. Thus, $\text{tr}(m_2)$ is a local coordinate at $\chi_n$, proving \textbf{(4)}. An integral basis vector for the kernel of the Jacobian is $(12, -(4n^3 + 12n^2 + 17n + 6), 4n^3 + 5n + 3)$. The Hessian of $\text{tr}([m_1, \ell_1])$ at $(-2, 2n+2, -2n)$ is given by 
\begin{equation}
	\begin{pmatrix}\frac{8n^2(n+1)^2(2n+1)^2}{9} & \frac{8n^2(n+1)(2n+1)}{3} & \frac{8n(n+1)^2(2n+1)}{3} \\ \frac{8n^2(n+1)(2n+1)}{3} & 8n^2 & 8n(n+1) \\ \frac{8n(n+1)^2(2n+1)}{3} & 8n(n+1) & 8(n+1)^2\end{pmatrix}
\end{equation}
Since this is nonzero on the tangent plane, it follows that the trace of $[m_1, \ell_1]$ is nonconstant at $\chi_n$, proving \textbf{(5)}. Now, since we already established that $\rho_n(\ell_1)$ and $\rho_n(\ell_2)$ are nonconjugate, it follows that $\rho_n|_{\langle m_1, \ell_1 \rangle}$ and $\rho_n|_{\langle m_2, \ell_2 \rangle}$ are nonconjugate representations in $SL_2(\mathbb{R})$, despite having the same character. Since $\text{tr}([m_1, \ell_1])$ is not constantly equal to 2 near $\chi_n$ in a small neighborhood of $\chi_n$ on $C_n$, $\rho_n|_{\langle m_i, \ell_i \rangle}$ for $i = 1, 2$ are irreducible with the same character. So these two restricted representations are conjugate in $SL_2(\mathbb{C})$, and they glue together to form the trace of a legitimate representation on the HNN extension gluing together $\langle m_1, \ell_1 \rangle$ to $\langle m_2, \ell_2 \rangle$, which is the fundamental group of $M_{-3, 3, 2n+1}$. Since $\chi_n$ cannot glue to form a representation, the nearby characters do not converge to a legitimate representation of $\pi_1(M_{-3, 3, 2n+1})$. Thus, they must converge to an ideal point, so $\chi_n$ corresponds to an ideal point on the character variety of $M_{-3, 3, 2n+1}$, and since nearby deformations have nonconstant longitude trace, this proves \textbf{(6)}. 
\end{proof}

\end{document}
